\documentclass[aspectratio=169]{beamer}

\usetheme{Madrid}
\usecolortheme{seahorse}

\usepackage[utf8]{inputenc}
\usepackage{amsmath, amssymb, amsfonts}
\usepackage{booktabs}
\usepackage{array}
\usepackage{tikz}
\usepackage{ragged2e}

% Accent color
\definecolor{AccentBlue}{RGB}{0,92,184}
\setbeamercolor{structure}{fg=AccentBlue}
\setbeamercolor{frametitle}{fg=AccentBlue}
\setbeamercolor{title}{fg=white,bg=AccentBlue}
\setbeamercolor{block title}{fg=white,bg=AccentBlue}
\setbeamercolor{block body}{bg=gray!5}

\setbeamertemplate{navigation symbols}{}

% Placeholder macro for figures/diagrams
\newcommand{\figplaceholder}[1]{%
  \begin{center}
    \begin{tikzpicture}
      \draw[dashed] (0,0) rectangle (10,5);
      \node[align=center] at (5,2.5) {\Large #1};
    \end{tikzpicture}
  \end{center}
}



\title[Housing Affordability]{Housing Affordability and Real Estate Market in Canada}
\subtitle{ECON 204 -- Contemporary Economic Issues}
\author{Muhebullah Karimzada}
%\institute{Winter 2026}
\date{Winter 2026}

\begin{document}

\maketitle

%================================================
\section{Introduction}

%------------------------------------------------
\begin{frame}{Introduction}
\transfade
\begin{itemize}[<+->]
  \item Housing affordability has become one of the most pressing economic challenges in Canada.
  \item Housing costs directly shape daily living standards, long-term financial security, and wealth accumulation.
  \item Housing is the largest single expenditure for most households and the dominant asset on household balance sheets.
  \item Changes in affordability affect inequality, intergenerational mobility, and financial stability.
  \item Since the mid-2000s, and especially after the COVID-19 pandemic, prices and rents rose faster than incomes in many regions.
  \item After 2022, higher interest rates slowed price growth, but affordability remained strained due to higher borrowing costs and tighter rental markets.
\end{itemize}
\end{frame}

%================================================
\section{What Is Housing Affordability?}

%------------------------------------------------
\begin{frame}{What Is Housing Affordability?}
\transfade
\begin{block}{Key Concept}
Housing affordability refers to the ability of households to obtain adequate housing without excessive financial strain, given their income, wealth, and access to credit.
\end{block}

\begin{itemize}[<+->]
  \item Affordability is a relationship, not a price level.
  \item High prices do not automatically imply unaffordability if incomes are high or financing costs are low.
  \item Moderate prices can be unaffordable if incomes are low or credit conditions are tight.
\end{itemize}
\end{frame}

%------------------------------------------------
\begin{frame}{Price-to-Income Measures}
\transfade
\begin{itemize}[<+->]
  \item A widely used indicator is the price-to-income ratio:
\end{itemize}

\vspace{0.2cm}
\[
PTI = \frac{\text{House Price}}{\text{Annual Household Income}}.
\]

\vspace{0.2cm}
\begin{itemize}[<+->]
  \item Historically, PTI ratios around 3--4 were common in many Canadian cities.
  \item Ratios above 7--8 reflect a breakdown in the traditional link between incomes and housing prices.
\end{itemize}
\end{frame}

%------------------------------------------------
\begin{frame}{Worked Numerical Example: Price-to-Income Ratio}
\transfade
\begin{block}{Worked Numerical Example}
Suppose the median house price in a Canadian city is \$900{,}000 and the median household income is \$110{,}000. The price-to-income ratio is:
\[
PTI = \frac{900{,}000}{110{,}000} \approx 8.18.
\]
This implies that the price of a typical home is more than eight times annual household income, a level widely regarded as severely unaffordable.
\end{block}
\end{frame}

%------------------------------------------------
\begin{frame}{Payment-Based Affordability}
\transfade
\begin{itemize}[<+->]
  \item Many households experience affordability through monthly cash flow rather than prices.
  \item Payment-based measures focus on the shelter-cost-to-income ratio:
\end{itemize}

\vspace{0.2cm}
\[
\text{Shelter Burden} = \frac{\text{Monthly Housing Costs}}{\text{Monthly Income}}.
\]

\vspace{0.2cm}
\begin{itemize}[<+->]
  \item This measure is especially important in Canada, where mortgage payments can change significantly at renewal.
\end{itemize}
\end{frame}

%------------------------------------------------
\begin{frame}{Worked Numerical Example: Shelter Burden}
\transfade
\begin{block}{Worked Numerical Example}
A household earns \$6{,}000 per month after tax. If mortgage payments, property taxes, and utilities sum to \$2{,}700 per month, the shelter burden is:
\[
\frac{2{,}700}{6{,}000} = 45\%.
\]
This household is housing-cost burdened even if the home value is high.
\end{block}
\end{frame}

%------------------------------------------------
\begin{frame}{Affordability Metrics: Summary Table}
\transfade
\begin{center}
\small
\setlength{\tabcolsep}{4pt}
\renewcommand{\arraystretch}{1.2}
\begin{tabular}{p{0.26\linewidth} p{0.34\linewidth} p{0.34\linewidth}}
\toprule
\textbf{Metric} & \textbf{Definition} & \textbf{Economic Interpretation} \\
\midrule
Price-to-income ratio (PTI) &
$PTI=\dfrac{\text{House Price}}{\text{Annual Household Income}}$ &
Captures the gap between prices and incomes; highlights long-run affordability stress. \\
\addlinespace[0.4em]
Shelter burden &
$\text{Burden}=\dfrac{\text{Monthly Housing Costs}}{\text{Monthly Income}}$ &
Captures cash-flow pressure; central when borrowing costs reset at mortgage renewal. \\
\bottomrule
\end{tabular}
\end{center}
\end{frame}

%================================================
\section{Housing as a Consumption Good and an Asset}

%------------------------------------------------
\begin{frame}{Housing as a Consumption Good and an Asset}
\transfade
\begin{itemize}[<+->]
  \item Housing plays two distinct economic roles at the same time.
  \item First, housing is a consumption good: households consume housing services (shelter, space, safety, location).
  \item If housing were only a consumption good, prices would be closely tied to fundamentals such as rents and incomes.
  \item Second, housing is also a store of wealth and an investment asset.
  \item Expectations about future house prices influence demand, similar to how expected stock returns influence stock demand.
\end{itemize}
\end{frame}

%------------------------------------------------
\begin{frame}{Why Prices Can Diverge from Rents or Incomes}
\transfade
\begin{itemize}[<+->]
  \item Because housing is both consumption and investment, demand depends on:
  \begin{itemize}[<+->]
    \item current ability to pay for shelter services, and
    \item expected future returns (expected appreciation).
  \end{itemize}
  \item When expected appreciation is high, the perceived long-run cost of ownership falls.
  \item Buyers may stretch financially today because they expect capital gains tomorrow.
  \item As a result, housing prices can diverge from fundamentals such as rents or incomes.
\end{itemize}
\end{frame}


%------------------------------------------------
\begin{frame}{Concrete Example: Expectations and Divergence}
\transfade
\begin{itemize}[<+->]
  \item Suppose two cities have identical rents and incomes.
  \item In City A, households expect house prices to be stable.
  \item In City B, households expect prices to rise by 5 percent per year.
  \item Demand for ownership will be higher in City B because buyers expect capital gains.
  \item House prices in City B will be higher and the price-to-rent ratio will be larger, even if rental fundamentals are identical.
\end{itemize}
\end{frame}

%================================================
\section{Demand-Side Drivers of Affordability in Canada}

%------------------------------------------------
\begin{frame}{Demand-Side Drivers of Housing Affordability in Canada}
\transfade
\begin{itemize}[<+->]
  \item Demand-side drivers increase households' willingness and ability to pay for housing at any given price.
  \item In Canada, demand forces have been unusually strong and persistent.
  \item The impact on affordability depends critically on the responsiveness of housing supply.
  \item When supply adjusts slowly, demand increases show up mainly as higher prices and rents rather than more housing.
\end{itemize}
\end{frame}

%------------------------------------------------
\begin{frame}{Population Growth and Household Formation}
\transfade
\begin{itemize}[<+->]
  \item Rapid population growth is a major driver (immigration, temporary residents, international students, natural increase).
  \item Population growth raises housing demand immediately because new households need housing upon arrival.
  \item Affordability deteriorates when household formation exceeds housing completions.
  \item Effects are pronounced in large metropolitan areas where inflows are concentrated.
\end{itemize}
\end{frame}

%------------------------------------------------
\begin{frame}{Income Growth and Income Dispersion}
\transfade
\begin{itemize}[<+->]
  \item Affordability depends on average income growth and the distribution of income.
  \item Higher-income households often have greater bidding power in desirable locations.
  \item Housing prices can reflect the purchasing power of the upper part of the income distribution rather than the median household.
  \item This can generate widespread affordability challenges even when aggregate income indicators appear strong.
\end{itemize}
\end{frame}

%------------------------------------------------
\begin{frame}{Credit Conditions and Interest Rates}
\transfade
\begin{itemize}[<+->]
  \item Housing affordability depends on income and how much households can borrow.
  \item Lower interest rates increase the mortgage size that can be serviced at a given payment, raising effective demand.
  \item Long periods of low rates increased borrowing capacity and allowed households to bid up prices.
  \item Rate increases reduce borrowing capacity, but can worsen affordability in the short run through renewal payment increases and tighter rental markets.
\end{itemize}
\end{frame}

%------------------------------------------------
\begin{frame}{Expectations of Future House Prices}
\transfade
\begin{itemize}[<+->]
  \item When households expect prices to rise, ownership becomes more attractive as both shelter and investment.
  \item Expected appreciation reduces the perceived long-run cost of ownership.
  \item This can generate self-reinforcing demand: rising prices attract buyers, which pushes prices higher.
  \item These dynamics can occur even among forward-looking and risk-aware households.
\end{itemize}
\end{frame}

%------------------------------------------------
\begin{frame}{Demographic and Life-Cycle Factors}
\transfade
\begin{itemize}[<+->]
  \item Life-cycle household formation patterns affect housing demand.
  \item Large cohorts entering prime home-buying ages increase demand for ownership housing.
  \item Longer life expectancy and aging in place reduce turnover in the existing housing stock.
  \item Affordability constraints delay homeownership for younger households, increasing rental demand and raising rents.
\end{itemize}
\end{frame}

%------------------------------------------------
\begin{frame}{Urban Concentration and Location Preferences}
\transfade
\begin{itemize}[<+->]
  \item Housing demand is increasingly concentrated in large urban centres where jobs and amenities cluster.
  \item Even if national supply expands, affordability can deteriorate locally when demand is spatially concentrated and supply responses are slow.
  \item Spillovers occur as households priced out of central cities relocate outward, transmitting demand pressure across regions.
\end{itemize}
\end{frame}

%================================================
\section{Supply Constraints and Price Dynamics}

%------------------------------------------------
\begin{frame}{Supply Constraints and Price Dynamics}
\transfade
\begin{itemize}[<+->]
  \item Supply constraints explain why affordability deteriorates sharply when demand increases.
  \item Housing supply is slow to adjust because housing is durable, land-intensive, and heavily regulated.
  \item In the short run, supply is relatively fixed, so prices rise to clear the market when demand increases.
  \item Supply elasticity determines whether adjustment occurs through quantities (construction) or prices.
\end{itemize}
\end{frame}

%------------------------------------------------
\begin{frame}{Housing Supply Elasticity}
\transfade
\begin{itemize}[<+->]
  \item Housing supply elasticity measures how responsive housing construction is to changes in prices.
  \item In elastic markets: higher prices lead to substantial increases in construction.
  \item In inelastic markets: higher prices generate little additional construction; adjustment occurs mostly through prices.
  \item Evidence indicates supply is relatively inelastic in many large Canadian metropolitan areas.
\end{itemize}
\end{frame}


%------------------------------------------------
\begin{frame}{Land-Use Regulation and Permitting}
\transfade
\begin{itemize}[<+->]
  \item Zoning rules that limit density, building heights, or reserve areas for single-family housing constrain construction.
  \item Lengthy approval processes and permitting uncertainty discourage rapid supply expansion.
  \item Developer delays reduce the ability to respond quickly to rising prices, reinforcing affordability pressures.
\end{itemize}
\end{frame}

%------------------------------------------------
\begin{frame}{Geographic Constraints and Construction Capacity}
\transfade
\begin{itemize}[<+->]
  \item Geography can limit developable land (coastlines, mountains, protected land).
  \item Expanding supply may require costly redevelopment or higher density, often facing political resistance.
  \item Construction capacity constraints (labour shortages, materials costs, infrastructure limits) can slow the pace of building.
  \item These constraints become binding during rapid demand growth.
\end{itemize}
\end{frame}

%------------------------------------------------
\begin{frame}{Persistence and Distributional Effects}
\transfade
\begin{itemize}[<+->]
  \item When supply responds slowly, prices do not quickly revert after a demand shock dissipates.
  \item Elevated prices can become embedded if households and investors expect constraints to persist.
  \item Scarcity rations access through prices: higher-income households absorb increases more easily.
  \item Lower- and middle-income households shift to smaller dwellings, longer commutes, or renting, increasing inequality and spatial segregation.
\end{itemize}
\end{frame}

%================================================
\section{Canadian Housing Market Case Studies}

%------------------------------------------------
\begin{frame}{Canadian Housing Market Case Studies}
\transfade
\begin{itemize}[<+->]
  \item Housing markets in Canada are highly regional; national averages conceal local differences.
  \item Comparing regions shows how shared forces (population growth, rates, income dynamics) interact with local supply elasticity and policy environments.
  \item There is no single Canadian housing market; there are distinct local markets responding to common pressures.
\end{itemize}
\end{frame}

%------------------------------------------------
\begin{frame}{Vancouver: Inelastic Supply and Price Capitalization}
\transfade
\begin{itemize}[<+->]
  \item Strong demand supported by high incomes, amenities, and sustained population growth.
  \item Supply constrained by geography (mountains and ocean), limited developable land, and restrictive land-use regulation.
  \item With inelastic supply, even small demand shocks can generate large price increases.
  \item Prices can become disconnected from local incomes due to investment demand and expectations of appreciation.
\end{itemize}
\end{frame}

%------------------------------------------------
\begin{frame}{Toronto and the Greater Golden Horseshoe: Growth Outpacing the Pipeline}
\transfade
\begin{itemize}[<+->]
  \item Rapid population growth driven by immigration and employment opportunities.
  \item Less severe geographic constraints than Vancouver, but supply limited by zoning, infrastructure coordination, and lengthy approvals.
  \item Large projects take years from approval to completion; starts may be strong while completions lag.
  \item Affordability can deteriorate if demand growth exceeds delivery, even when construction expands.
\end{itemize}
\end{frame}

%------------------------------------------------
\begin{frame}{Montreal: From Relative Affordability to Tightening Markets}
\transfade
\begin{itemize}[<+->]
  \item Historically relatively affordable, especially in rentals, due to large rental stock and slower population growth.
  \item Recent pressures intensified with faster population growth and prolonged low rates.
  \item Higher rates after 2022 moderated price growth, but affordability challenges persisted via stronger rental demand and higher renewal payments.
  \item Relative affordability advantages are not permanent; conditions can shift quickly when supply is slow.
\end{itemize}
\end{frame}

%------------------------------------------------
\begin{frame}{Prairie Cities: Elastic Supply with Emerging Pressures}
\transfade
\begin{itemize}[<+->]
  \item Calgary and Edmonton historically more affordable due to more elastic supply (abundant land, fewer restrictions, faster permitting).
  \item Recently, strong migration and population growth increased demand sharply.
  \item Supply responded more than in constrained markets, but not enough to prevent rising prices and rents.
  \item Even in relatively elastic markets, affordability can deteriorate when demand accelerates rapidly.
\end{itemize}
\end{frame}

%------------------------------------------------
\begin{frame}{Atlantic Canada: Small Markets and Large Shocks}
\transfade
\begin{itemize}[<+->]
  \item Smaller markets historically had modest price growth.
  \item Pandemic-era remote work, migration, and demographic shifts raised demand.
  \item Supply was slow to adjust due to limited construction capacity and smaller pipelines.
  \item Lower initial price levels made percentage price increases appear especially large.
\end{itemize}
\end{frame}

%------------------------------------------------
\begin{frame}{Case Studies: Comparison Table}
\transfade
\begin{center}
\scriptsize
\setlength{\tabcolsep}{3pt}
\renewcommand{\arraystretch}{1.2}
\begin{tabular}{p{0.18\linewidth} p{0.26\linewidth} p{0.26\linewidth} p{0.24\linewidth}}
\toprule
\textbf{Region} & \textbf{Core Demand Forces} & \textbf{Key Supply Constraints} & \textbf{Affordability Dynamics} \\
\midrule
Vancouver & Amenities, incomes, population growth, investment demand & Geography + restrictive density constraints & Small demand shocks lead to large price increases; persistent high PTI \\
\addlinespace[0.35em]
Toronto GGH & Immigration + jobs, sustained inflows & Approvals, zoning, long project timelines & Pipeline lags demand; prices and rents remain under pressure \\
\addlinespace[0.35em]
Montreal & Accelerating demand + low-rate legacy & Slow adjustment in supply; tighter rentals & Shift from relative affordability to strained markets \\
\addlinespace[0.35em]
Prairie Cities & Migration-driven demand increases & Capacity constraints when demand surges & More quantity response, but still rising prices/rents recently \\
\addlinespace[0.35em]
Atlantic Canada & Remote-work and migration shocks & Limited construction capacity & Outsized percentage price increases in smaller markets \\
\bottomrule
\end{tabular}
\end{center}
\end{frame}

%================================================
\section{Monetary Policy and Mortgage Renewal}

%------------------------------------------------
\begin{frame}{Monetary Policy and Mortgage Renewal}
\transfade
\begin{itemize}[<+->]
  \item Monetary policy affects housing affordability through several channels.
  \item The mortgage renewal channel is particularly important in Canada.
  \item This channel helps explain why affordability can worsen even if house prices stop rising or fall.
  \item Canadian mortgages expose households to interest rate risk at renewal rather than locking costs for the full loan life.
\end{itemize}
\end{frame}

%------------------------------------------------
\begin{frame}{Renewal Mechanics in Canada}
\transfade
\begin{itemize}[<+->]
  \item Many mortgages have fixed rates for short terms (commonly five years), while amortizations are much longer (often 25 or 30 years).
  \item Households periodically renew at prevailing market rates.
  \item Rate hikes do not hit all borrowers immediately; impacts arrive as mortgages renew.
  \item Renewal at higher rates can raise monthly payments substantially even if balances declined modestly.
\end{itemize}
\end{frame}

%------------------------------------------------
\begin{frame}{Price Channel vs. Payment Channel}
\transfade
\begin{itemize}[<+->]
  \item Higher rates reduce borrowing capacity for new buyers, slowing demand and price growth (price channel).
  \item For existing borrowers, higher rates raise required payments at renewal (payment channel).
  \item Prices and affordability can move in opposite directions:
  \begin{itemize}[<+->]
    \item prices stabilize or decline, while
    \item affordability worsens as monthly payments rise.
  \end{itemize}
\end{itemize}
\end{frame}



%------------------------------------------------
\begin{frame}{Distributional Consequences of Renewal Risk}
\transfade
\begin{itemize}[<+->]
  \item Recently purchased households and highly leveraged borrowers are more exposed to renewal risk.
  \item These households tend to be younger and more vulnerable to payment increases.
  \item Long-time homeowners with small balances or no mortgage are largely insulated.
  \item Monetary tightening can redistribute stress and exacerbate intergenerational inequality.
\end{itemize}
\end{frame}

%------------------------------------------------
\begin{frame}{Renters and the Renewal Channel}
\transfade
\begin{itemize}[<+->]
  \item Higher interest rates raise financing costs for landlords.
  \item In tight rental markets, some higher costs may pass through into higher rents.
  \item Higher rates can delay new construction, constraining future rental supply and reinforcing rent pressure.
  \item Effects unfold with long and uneven lags; stress can peak well after a tightening cycle begins.
\end{itemize}
\end{frame}

%================================================
\section{International Housing Affordability Crises}

%------------------------------------------------
\begin{frame}{International Housing Affordability Crises}
\transfade
\begin{itemize}[<+->]
  \item The international housing affordability crisis refers to widespread and simultaneous deterioration in affordability across many advanced economies over the past two decades.
  \item Affordability problems have emerged across countries with different institutions and policy frameworks.
  \item This suggests common economic forces rather than isolated national policy failures.
  \item A core feature is divergence: housing costs rise faster than household incomes.
\end{itemize}
\end{frame}

%------------------------------------------------
\begin{frame}{A Common Global Pattern}
\transfade
\begin{itemize}[<+->]
  \item Demand concentrates in major urban centres with jobs and amenities.
  \item Supply responds slowly due to zoning, regulation, infrastructure limits, and construction bottlenecks.
  \item Long periods of low interest rates increased borrowing capacity and raised asset prices.
  \item With inelastic supply, demand increases translate mainly into higher prices rather than more housing.
\end{itemize}
\end{frame}

%------------------------------------------------
\begin{frame}{The United States: Regional Divergence}
\transfade
\begin{itemize}[<+->]
  \item In some cities, supply is highly constrained by zoning, community opposition, and lengthy approvals.
  \item In these markets, rising demand leads mainly to higher prices and severe affordability pressure.
  \item In more flexible markets, price increases have been smaller.
  \item The contrast highlights the central role of supply elasticity.
\end{itemize}
\end{frame}

%------------------------------------------------
\begin{frame}{United Kingdom: Rental Affordability and Generational Inequality}
\transfade
\begin{itemize}[<+->]
  \item Rising house prices push more younger households into renting.
  \item Rental demand rises faster than rental supply, raising rents sharply in major areas.
  \item Renters can face higher effective inflation than homeowners.
  \item The result is intensified generational inequality in housing cost burdens.
\end{itemize}
\end{frame}

%------------------------------------------------
\begin{frame}{Australia and New Zealand: Extreme Price-to-Income Ratios}
\transfade
\begin{itemize}[<+->]
  \item Cities reach exceptionally high price-to-income ratios.
  \item Drivers include strong population growth, urban concentration, constrained supply, and financial conditions that favor housing as an investment asset.
  \item These cases show that affordability crises can persist even in high-income countries with stable institutions.
\end{itemize}
\end{frame}

%------------------------------------------------
\begin{frame}{Europe: Different Institutions, Similar Outcomes}
\transfade
\begin{itemize}[<+->]
  \item European countries vary in tenure systems, yet affordability problems have emerged across many cities.
  \item In major metropolitan areas, limited supply combined with strong demand pushed prices and rents above income growth.
  \item Convergence across different systems suggests interaction of global demand pressures with local supply constraints.
\end{itemize}
\end{frame}

%------------------------------------------------
\begin{frame}{Why Affordability Has Become Structural}
\transfade
\begin{itemize}[<+->]
  \item The current episode is often described as structural rather than purely cyclical.
  \item Affordability worsened persistently over time.
  \item Younger and lower-income households are increasingly excluded from homeownership.
  \item Housing costs reshape labour mobility, family formation, and inequality.
  \item In many cases, affordability does not return to pre-boom levels even after markets cool.
\end{itemize}
\end{frame}

%------------------------------------------------
\begin{frame}{Implications for Canada}
\transfade
\begin{itemize}[<+->]
  \item Canada fits the international pattern: strong demand concentrated in large cities, constrained supply, and financial channels amplifying price movements.
  \item An international lens clarifies that the issue is broader than a single-country policy failure.
\end{itemize}
\end{frame}

%================================================
\section{Conclusion}

%------------------------------------------------
\begin{frame}{Conclusion}
\transfade
\begin{itemize}[<+->]
  \item Housing affordability in Canada reflects interaction of strong demand growth, constrained supply, and powerful financial channels.
  \item Prices, rents, and payments respond differently to economic shocks.
  \item Affordability can deteriorate even when some indicators improve.
  \item Understanding these mechanisms is essential for evaluating policy responses and interpreting housing market developments in Canada and abroad.
\end{itemize}
\end{frame}



\end{document}
